\documentclass[11pt]{article}
\usepackage[a4paper,margin=1in]{geometry}
\usepackage{amsmath,amssymb,amsthm}
\usepackage[T1]{fontenc}
\usepackage{lmodern}
\usepackage[hidelinks]{hyperref}
\setlength{\emergencystretch}{3em}

\newtheorem{proposition}{Proposition}
\newtheorem{lemma}[proposition]{Lemma}

\title{A Support-Endpoint Counterexample to Conjecture 00000007737}
\author{Cosica}
\date{2026-10-08}

\begin{document}
\maketitle

\begin{abstract}
Conjecture 00000007737 asserts a formula for the lower endpoint $x_k$ of the
support of the $k$-fold free multiplicative convolution of the uniform law on
$[0,1]$. At $k=2$ its formula gives $x_2=16e^{-2}>1$. This is impossible:
the free multiplicative convolution of two probability measures supported on
$[0,1]$ is itself supported on $[0,1]$. Thus its lower endpoint is at most
$1$. This report gives the counterexample and an accompanying Lean~4
formalization of the endpoint contradiction.
\end{abstract}

\section{The asserted formula}

The conjecture defines $\mu_k=\nu\boxtimes\cdots\boxtimes\nu$, where $\nu$
is the uniform probability law on $[0,1]$, and asserts
\[
  x_k=\exp(-2k+2) k^{2k/(k-1)}.
\]
There is no stated restriction excluding $k=2$. Substitution gives
\begin{equation}\label{eq:value}
  x_2=\exp(-2)2^4=\frac{16}{e^2}.
\end{equation}

\section{The support obstruction}

\begin{lemma}\label{lem:contraction}
If $a$ and $b$ are positive contractions in a $C^*$-algebra, then
$a^{1/2}ba^{1/2}$ is a positive contraction.
\end{lemma}

\begin{proof}
Positivity follows from
\[
  a^{1/2}ba^{1/2}=(b^{1/2}a^{1/2})^*(b^{1/2}a^{1/2}).
\]
Since $0\le b\le1$, conjugation by $a^{1/2}$ gives
\[
  0\le a^{1/2}ba^{1/2}\le a^{1/2}1a^{1/2}=a\le1.
\]
\end{proof}

\begin{proposition}\label{prop:support}
The support of $\mu_2=\nu\boxtimes\nu$ is contained in $[0,1]$.
Consequently its lower endpoint is at most $1$.
\end{proposition}

\begin{proof}
Use the operator definition of free multiplicative convolution
(see, for example, the introduction of~\cite{Ji}): take two copies of
the coordinate function $t\mapsto t$ in $C([0,1])$, equipped with integration
against $\nu$, in their reduced free product probability space. The images
$a,b$ are free positive contractions with distribution $\nu$; positivity and
the inequality $a,b\le1$ are preserved by the unital star homomorphisms.
The distribution of
$a^{1/2}ba^{1/2}$ is $\nu\boxtimes\nu$. Lemma~\ref{lem:contraction} places
the spectrum, hence the support of this distribution, in $[0,1]$. A nonempty
subset of $[0,1]$ has infimum at most $1$.
\end{proof}

\section{Numerical contradiction}

\begin{lemma}\label{lem:numerical}
\[
  \frac{16}{e^2}>1.
\]
\end{lemma}

\begin{proof}
The elementary estimate $e<3$ yields $e^2<9<16$. Dividing the strict
inequality $e^2<16$ by the positive number $e^2$ proves the claim.
\end{proof}

Combining Proposition~\ref{prop:support}, Lemma~\ref{lem:numerical}, and
\eqref{eq:value} gives $x_2\le1<x_2$, a contradiction. Hence the claimed
formula, and therefore Conjecture 00000007737 as a conjunction containing
that formula, is false.

\section{Lean 4 verification}

The included Lean project uses Lean~4.33.1 and Mathlib~4.33.1.
It proves a stronger operator statement: for every pair of positive
contractions in any nontrivial unital complex $C^*$-algebra, every scalar
spectral probability law of their positive product has lower support endpoint
different from the conjectured value at $k=2$.
Here a scalar spectral law is represented as a probability measure $\rho$
on the real spectrum of $c=a^{1/2}ba^{1/2}$, pushed forward along the
inclusion into $\mathbb R$. This is the standard representation of an
operator's distribution in a state, obtained by restricting the state to
the commutative algebra generated by $c$.
The formal proof establishes:
\begin{itemize}
  \item \texttt{positiveProduct\_nonneg} and
        \texttt{positiveProduct\_le\_one} prove $0\le c\le1$;
  \item \texttt{positiveProduct\_spectrum\_subset} proves
        $\sigma_{\mathbb R}(c)\subseteq[0,1]$;
  \item \texttt{spectralLaw\_support\_subset} and
        \texttt{spectralLaw\_support\_nonempty} prove the corresponding
        support bound and nonemptiness for the probability law;
  \item \texttt{claimedEndpoint\_gt\_one} proves $1<16/\exp(2)$, using
        Mathlib's certified bound $\exp(1)<3$;
  \item \texttt{conjecturedEndpoint\_two} checks the substitution into the
        original exponential and real-power formula;
  \item \texttt{conjecture\_00000007737\_endpoint\_false} combines these
        facts for every such operator pair and spectral probability law.
\end{itemize}
The construction of free products and identification of their spectral laws
with free multiplicative convolution are used at the level of the standard
definition in the report; the Lean project does not implement a separate
free-convolution operation. Its universal operator theorem needs neither
freeness nor uniformity, so it applies to the defining realization above.
In particular, the unit-support bound is proved, rather than assumed as a
hypothesis of the final operator theorem.

The project introduces no additional axioms and uses no \texttt{sorry},
\texttt{admit}, or \texttt{native\_decide}. The printed axiom dependencies
are exactly \texttt{propext}, \texttt{Classical.choice}, and
\texttt{Quot.sound}. Run \texttt{lake build} to check the proof and print
these dependencies.

\begin{thebibliography}{1}
\bibitem{Ji} H. C. Ji, \emph{Regularity Properties of Free Multiplicative
Convolution on the Positive Line}, International Mathematics Research
Notices 2021(6), 4522--4563. Preprint:
\url{https://arxiv.org/abs/1903.02326}.
\end{thebibliography}

\end{document}
